\chapter{The Field Closes}

\section{The second variation}

Part V opens on the field --- and opens on it lightly, because the field, read in full, is the physical volume's
climax, not this one's. What this volume has to say about the field is small and structural, and it can be said
quickly. The residue the book has read all along is a first difference; the field is its \emph{second} variation ---
the curvature of the residue, the residue read one order deeper. The field is not a new object dropped into the
construction at the end. It is the same residue, differentiated once more.

Take that plainly, because it is the whole of what this volume claims the field to be. A difference has a rate; the
rate has a rate; and the second of those --- how the residue's own change is changing --- is the second variation. The
physical gauge, reading this second variation as a quantity in the world, calls it a field. But underneath the name it
is exactly what the construction has been carrying: the residue, one derivative deeper. Nothing about the machinery
changes to produce it; the field is a reading of the same leftover, taken at second order.

And here is the one load-bearing thing this volume keeps, and it is a claim about \emph{finiteness}. You might expect a
field to be a continuum --- a smooth quantity spread over space, to be integrated. It is not, at this tier. The
construction renders the second variation as a small, fixed number of \emph{channels} --- and the number is not the
book's choice but the code's: the construction lists exactly four, names them, and stops. Each channel is a ratio of
the kind the last chapter built --- a floor and a remainder, a magnitude with an orientation --- so the field, here, is
four ratios: four finite readings, held as four remainders, and no continuum anywhere in sight.

This is worth pausing on against something the book said earlier. In the chapter on the faces, the residue was read
four ways, and the book was careful to mark that ``four'' as its own --- the count of the gauges, an expository choice,
not a number the machine committed to. Here the four is different. These four channels are enumerated in the
construction itself; the code names them and decides between them. So the earlier four faces --- the book's four ---
close, at the field, into four channels that are the machine's four. It is a satisfying seam, but the distinction must
be kept exactly: the faces were a reading given four names by the book; the channels are four things the code actually
holds. The expository four has found a code-real four to land on, and that is a genuine confirmation --- but it does not
reach back and make the faces themselves proved. Two different fours, meeting honestly at the field.

What the four channels \emph{mean} as a field --- the electromagnetism the physical gauge reads in them, the continuous
quantity it integrates, the coupling it extracts --- is the next volume's climax, and this one names the second
variation, names its four finite channels, and hands the field forward. What stays here is only the finiteness: the
field, at this tier, is the second variation kept \emph{finite} --- four channels, four ratios, not a smooth continuum
--- and that finiteness is the same discipline the whole book has run on. When the far chapters read the coupling, they
will read it off this finite field, at the machine's own floor, and never off a continuum. This volume keeps that fact
and no more.

In the two verbs the field is a fourfold funge of the residue's second variation, each channel a bag with a remainder
--- the residue tanged out of each, exactly as the apparatus tanged it in the last chapter. The continuous field the
physical volume will read is that volume's funge to make; this one keeps the four finite bags and the residue in each.
So the field, opening Part V, arrives not as a smooth infinity but as four finite readings of the residue's curvature
--- the same leftover, one order deeper, held four ways and no way rounded. The field closes, in this volume, by
staying finite.

\section{The dual side present}

The last section gave the field its electric side --- the four finite channels of the second variation. But a field has
two sides, and the construction carries both. This short section names the other one --- the magnetic side --- shows
that it, too, is kept finite, and then names the place where the two sides meet, which is the number the whole
instrument has been built to reach. It names that meeting, and, as it has done from the first page, it declines to give
the number.

The magnetic side is read the way the book reads everything: as a ratio. The physical setup behind it is a Meissner
arrangement --- a surface current that expels an applied field --- and the physics of that expulsion, the
superconducting mechanism, is the next volume's to explain. What this volume keeps is the arithmetic left when the
physics is set aside: a ratio, the applied field over the shielding quantity, a floor and a remainder like every
reading before it. The magnetic side is not a new kind of thing; it is the same apparatus, pointed at the field's other
half.

And so both sides of the field are finite. The electric side was four ratios; the magnetic side is a ratio; and the
whole field, taken together, is held as a handful of finite readings --- never as the continuous field the physical
gauge will one day integrate. This is the same discipline the book has kept throughout, now applied to the last object
it had left to apply it to: the field, both its sides, stays finite on this bench. There is no continuum here to reach
into; there are ratios, and their remainders, and nothing smooth between them.

Now the place the whole book has been walking toward, and the discipline that has governed every step of the walk.
Where the electric side and the magnetic side meet --- where the two finite halves of the field are set against each
other --- is the electromagnetic coupling. That coupling is the number the instrument was built to read, and the
construction, at this point, has everything it needs to ask for it: both sides, both finite, both in hand. And here the
book does exactly what it has done since its opening pages. It names the coupling --- the meeting of the two sides ---
and it carries not one digit of its value. The magnitude is the next chapter's, and it is held blind here as strictly
as it has been held everywhere: the coupling is named, its two finite sides are in hand, and its number is left
un-written, waiting. The machine can now ask; the asking is one chapter away; and even when it asks, the answer will
come back not as a point but as a bracket. What this section delivers is the question, fully assembled, and nothing of
the answer.

In the two verbs the magnetic side funges into its ratio, the residue tanged out of it exactly as the apparatus tanged
the electric residue out. And the coupling itself is a funge the machine deliberately leaves open --- the two sides
bagged toward a single number the machine will not, here, complete to a value. It is the same refusal the bracket has
always been: a bag the machine holds without pinning, a number named without being fixed. So Part V's opening chapter
closes with the field whole and finite, its two sides in hand, and the coupling named and un-numbered between them ---
everything assembled to read the number, and the number, honestly, still withheld.

\section{Asking, not delivering}

The last section left the field whole: both sides finite, both in hand, and the coupling named between them without a
number. This closing section performs the one act the assembled field is for --- it \emph{asks} for the coupling ---
and then does the thing the whole book has trained the reader to expect. It does not answer. The field's closure poses
the question, fully and with everything in place, and hands over the question, not the magnitude.

The closure is a \emph{contraction}. The machine gathers the four channels of the electric side and the ratio of the
magnetic side and draws them together toward a single dimensionless number --- the coupling. The construction carries
an operation that does exactly this: it takes the assembled field and contracts it toward the one number the two sides
determine between them. And that gathering-and-pointing \emph{is} the asking. To assemble the field and contract it
toward the coupling is to put the question --- what is the coupling? --- in the only form the machine can put it: an
apparatus, built and aimed.

And here is the discipline that keeps the asking honest, and it is the anti-crank heart of the whole method. The
construction is explicit that the calibration behind the contraction is generated \emph{before} the coupling is read,
not fitted \emph{after} it. The machine does not build its apparatus around a number it already wanted and then present
that number as a discovery; it assembles the apparatus first, from the two finite sides, honestly, and only then asks
what the assembly yields. That order is the whole difference between an instrument and a fake. A fake tunes its
apparatus to the answer it means to announce, so that the announcement is guaranteed; this machine builds its apparatus
before it looks, and takes whatever the looking gives. Build, then ask --- never ask, then build to match.

And having asked, honestly, the closure does not deliver. The coupling's value is the next chapter's, held blind here
as strictly as it has been held on every page before this one. The field's closure poses the question --- everything
assembled, the apparatus aimed --- and stops there, handing the question forward and keeping no number for itself. Nor
will the next chapter, when it answers, hand back a point: the answer will be a bracket, floored at the machine's own
resolution and reruning to the same interval every time it is asked --- the same honest object the book built parts
ago, a lower, an upper, and no obligation to name anything between. The magnitude the closure withholds is withheld not
because the machine is coy but because the honest answer, when it comes, is not a magnitude. It is a bound.

In the two verbs the contraction \emph{funges} the whole field into a single question --- the two sides, the four
channels and the one ratio, bagged together as ``the coupling.'' And the machine, having funged them into the question,
declines the last step: it will not \emph{tange} the question down to a value. Asking is the honest funge; delivering a
magnitude would be the over-tange the book has refused since it first chose the bag over the pin. So Part V's opening
chapter ends with the instrument complete and the question posed --- the field contracted to the coupling, the coupling
held open, and the number, one chapter from here, still un-taken. Everything to ask with, and nothing yet answered.
